\begingroup
\setlength{\emergencystretch}{2em}
\subsection*{Архитектура и организация компьютера}
\label{lecture04-journal}
\textit{Сквозной пример: разработка журнала оценок группы.}
Напишем программу без гостевой ОС. Исполняемый код --- 16-битный NASM
в real mode на эмулированном QEMU PC (i440FX); C встречается лишь как
\emph{пояснение алгоритма}. Сохраняем действующий текст основной лекции №4;
здесь те же архитектурные понятия открываются при разработке шести
последовательных версий журнала. Студенты обозначены \texttt{S01}--\texttt{S16},
у каждого четыре работы и оценка 0--10 либо «не выставлена».

\begin{center}\small
\begin{tabular}{|p{0.23\linewidth}|p{0.59\linewidth}|}\hline
Минуты & Потребность и теория \\ \hline
0--7 & Задача и простая аппаратная модель \\
7--20 & Запуск: BIOS, сектор, CS:IP и первый символ \\
20--35 & Таблица: байты, адреса и текстовая видеопамять \\
35--48 & Ввод через BIOS и разбор скан-кода \\
48--60 & Собственный IRQ1 и очередь событий \\
60--72 & Средние, регистры, стек, код и кэш \\
72--86 & Долговременный сектор, BIOS-диск и проверка данных \\
86--90 & Итоги и связь с операционной системой \\ \hline
\end{tabular}
\end{center}

\paragraph{Контракт учебной машины.}
У QEMU PC 16 MiB RAM, BIOS, VGA-текстовый экран, PS/2-клавиатура,
PIC, COM1 для диагностики, загрузочная дискета 1.44 MiB и отдельный
IDE-совместимый диск 1 MiB для оценок. Настоящее устройство компьютера
сложнее, но выбранный набор позволяет увидеть, кто исполняет инструкции,
где лежат данные и как взаимодействуют компоненты. BIOS --- прошивка,
не гостевая ОС. Для тестов запускаем \texttt{qemu-system-i386}; название
команды не исключает исполнения 16-битных инструкций в real mode.
Ни одну характеристику производительности реального SSD или кэша не
измеряем по скорости QEMU.

\paragraph{Версия 1: начать исполнение и вывести одну оценку.}
После сброса CPU выбирает инструкции BIOS. BIOS выполняет базовую
инициализацию, выбирает дискету согласно порядку загрузки, размещает
её первый сектор по адресу \texttt{0000:7C00} и передаёт туда управление.
Его последние два байта имеют сигнатуру \texttt{55 AA}; сама программа
слишком велика для сектора 512 байтов. Поэтому наш загрузочный сектор
\texttt{boot16.asm} читает сектора 2--18 той же дорожки в RAM через
BIOS \texttt{INT 13h}, проверяет результат и делает дальний переход
к \texttt{1000:0000}:
\begin{lstlisting}[basicstyle=\ttfamily\small]
mov ax, 0x1000
mov es, ax
xor bx, bx
mov ah, 0x02
mov al, 17
mov cl, 2
int 0x13
jc boot_error
jmp 0x1000:0x0000
\end{lstlisting}
В полном листинге дополнительно задаются цилиндр, головка и номер
загрузочной дискеты \texttt{DL}, а также проверяется число прочитанных
секторов. За пределы первой дорожки не выходим; сборка отвергает программу
длиннее $17\cdot512=8704$ байтов. \textbf{Наш код запускается без GRUB},
хотя BIOS сам по себе остаётся частью учебной машины.

Адрес в real mode рассчитывается как $16\cdot\text{segment}+\text{offset}$.
Например, \texttt{1000:0000} обозначает физический \texttt{0x10000}.
Инструкции адресует \texttt{CS:IP}; \texttt{DS} указывает данные программы,
\texttt{SS:SP} --- её стек. В начале основного NASM-кода \texttt{DS=CS},
\texttt{SS=7000h}, \texttt{SP=F000h}. Глобальные данные и стек поэтому
не обязаны занимать одно место. Это учебная конфигурация, не модель
адресного пространства процесса современной ОС.

Чтобы вывести \texttt{S01  grade 07}, сначала используем готовый сервис
прошивки. В процедуре \texttt{bios\_text} символы читаются из строки
\texttt{DS:SI} и передаются в BIOS:
\begin{lstlisting}[basicstyle=\ttfamily\small]
lodsb                    ; AL = next character
test al, al
jz .done
mov ah, 0x0e             ; BIOS teletype
mov bx, 0x0007
int 0x10
\end{lstlisting}
\texttt{INT 10h} --- \emph{программная} инструкция обращения к BIOS;
это не аппаратное уведомление о событии. Цифра на экране --- ASCII-символ
\texttt{'7'}, а оценка 7 станет числовым значением данных лишь на следующей
стадии. Диагностический COM1 отдельно выдаёт \texttt{BOOT STAGE=1} и
\texttt{READY}, чтобы тест наблюдал действительный запуск без распознавания
изображения.

\paragraph{Версия 2: массив оценок и прямой вывод.}
Массив содержит 16 строк по 4 однобайтовых поля: $16\cdot4=64$ байта.
Коды 0--10 означают выставленные оценки, а \texttt{FF} означает
«не выставлена»; нуль \emph{не} может служить признаком пропуска.
Смещение выбранного элемента от начала массива равно $4i+j$.
В NASM это резервируется и инициализируется как
\texttt{grades: times 64 db 0xff}. Пояснение на C:
\begin{lstlisting}[language=C,basicstyle=\ttfamily\small]
uint8_t grade[16][4]; /* algorithm only */
\end{lstlisting}
Учебное значение \texttt{grade[1][2]} имеет смещение 6, но полный адрес
определяется сегментом \texttt{DS} и реальным размещением метки
\texttt{grades}. Адрес, байт значения и вывод на экран --- три разные вещи.

После знакомства с VGA текстовый вывод делаем самостоятельно: буфер
устройства начинается с физического \texttt{0xB8000}, то есть \texttt{B800:0000}.
На одну позицию приходятся код символа и байт атрибута. Строка экрана имеет
$80\cdot2=160$ байтов; адрес позиции внутри видеосегмента равен
$160\cdot\text{row}+2\cdot\text{col}$.
\begin{lstlisting}[basicstyle=\ttfamily\small]
mov ax, 0xb800
mov es, ax
; DI already contains the cell offset
lodsb                    ; AL = character at DS:SI
mov ah, 0x0f             ; color attribute
stosw                    ; AX -> ES:DI
\end{lstlisting}
Это ключевые строки реальной функции \texttt{vga\_text}.
Процедура \texttt{redraw} перебирает оценки в RAM и создаёт
\emph{отображение} в VGA; дисплей не становится долговременным хранилищем.
Для современной видеокарты потребовались бы иные драйверы и, возможно,
отдельная VRAM.

\paragraph{Версия 3: изменяем оценку с клавиатуры.}
Через BIOS \texttt{INT 16h} сначала спрашиваем о готовом событии
(\texttt{AH=01h}), затем читаем его (\texttt{AH=00h}):
\begin{lstlisting}[basicstyle=\ttfamily\small]
mov ah, 0x01
int 0x16
jz .loop
xor ah, ah
int 0x16
mov al, ah              ; AL = scan code for handle_scan
call handle_scan
\end{lstlisting}
PS/2-контроллер сообщает BIOS о клавише; BIOS предоставляет программе
высокоуровневый сервис. Клавиша \texttt{A} даёт оценку 10, цифры 0--9
задают соответствующее значение, Backspace записывает \texttt{FF},
стрелки выбирают работу и студента. Наш \texttt{handle\_scan} переводит
скан-код в действие. ASCII, скан-код и числовая оценка не тождественны.
Контроллер и CPU обмениваются через аппаратный интерфейс с адресом/портом,
данными и управляющей информацией; современные последовательные соединения
могут упаковывать это в сообщения по одним линиям. BIOS скрывал порты
\texttt{64h} и \texttt{60h} от первых версий программы.

\paragraph{Версия 4: собственное аппаратное прерывание.}
Теперь обработку PS/2 IRQ1 берёт наш код, поэтому BIOS \texttt{INT 16h}
больше не используется: его буфер клавиш уже не пополняется стандартным
обработчиком. В real mode таблица векторов \texttt{IVT} начинается с
\texttt{0000:0000}. Четыре байта по смещению $9\cdot4=36$ задают IP и CS
обработчика IRQ1 в совместимом PIC. При изменении IVT прерывания временно
выключаем через \texttt{cli}, затем включаем через \texttt{sti}.
После сохранения регистров короткий обработчик считывает скан-код
из порта \texttt{60h}, кладёт его в кольцевую очередь и подтверждает
PIC (EOI); возврат выполняет \texttt{iret}. Упрощённая схема кода:
\begin{lstlisting}[basicstyle=\ttfamily\small]
irq1:
    push ax
    push ds
    mov ax, cs
    mov ds, ax
    in al, 0x60
    ; enqueue the byte when there is room
    mov al, 0x20
    out 0x20, al
    pop ds
    pop ax
    iret
\end{lstlisting}
Реальный исходник также проверяет бит готовности, сохраняет \texttt{BX}
и содержит полный код очереди. Обработчик не перерисовывает экран:
основной цикл вызывает \texttt{dequeue}, затем \texttt{handle\_scan}.
Буфер из 64 байтов хранит не более 63 событий, потому что одна позиция
разделяет состояния «пусто» и «полно». Основной цикл всё ещё опрашивает
\emph{очередь}, но не PS/2-порт. Внешний IRQ1, намеренное \texttt{INT 16h}
и исключение от CPU имеют разные причины, хотя используют механизм
передачи управления.

\paragraph{Версия 5: статистика, регистры и память.}
Среднее вычисляем только из выставленных оценок. Если количество $n=0$,
показываем \texttt{--.--}, иначе в сотых:
\[
\left\lfloor\frac{100\cdot\mathrm{sum}+\lfloor n/2\rfloor}{n}\right\rfloor.
\]
Например, 0 и 10 дают 5.00; 0, 1 и 1 дают 0.67. Такое целочисленное
вычисление не требует библиотеки float. В NASM четыре поля читаются
из \texttt{DS:grades+BX}, байт \texttt{FF} пропускается, число 0
учитывается. Цикл из диагностики COM1:
\begin{lstlisting}[basicstyle=\ttfamily\small]
mov cx, WORKS
xor dx, dx             ; DH=sum, DL=count
.sum:
    mov al, [grades+bx]
    cmp al, UNSET
    je .skip
    inc dl
    add dh, al
.skip:
    inc bx
    loop .sum
\end{lstlisting}
Регистры держат текущую сумму; \texttt{CS:IP} выбирает байты машинных
инструкций, а \texttt{DS:grades+BX} --- байты данных. В строгой гарвардской
машине эти адресные пространства были бы разделены, в выбранном x86 PC
код и данные разделяют обычную RAM. Современные CPU могут иметь разные
L1I и L1D; QEMU-тест проверяет результат, но не измеряет настоящие
кэш-промахи. Последовательные 64 байта оценок демонстрируют локальность:
при условной 64-байтовой кэш-линии массив может занимать одну линию,
если выровнен, иначе две.

Стек используется вызовами, прерываниями и сохранением регистров.
В \texttt{entry} заранее устанавливаются \texttt{SS:SP}; данные оценок
лежат в отдельном статическом резерве. Для фиксированной группы
аллокатор кучи не требуется. Если разрешить динамическое число студентов,
понадобилось бы управление свободными блоками и временем жизни.

\paragraph{Версия 6: долговременное хранение.}
RAM энергозависима, поэтому оценки пишутся в отдельный диск.
Загрузочная программа на виртуальной дискете (диск BIOS 00h), оценки
в секторе LBA 1 первого HDD (диск BIOS 80h). BIOS \texttt{INT 13h}
с расширенным адресным пакетом DAP предоставляет чтение \texttt{AH=42h}
и запись \texttt{AH=43h}; флаг CF сообщает об ошибке. DAP задаёт
число секторов, физический адрес буфера \texttt{1000:sector} и номер LBA.
Здесь нет файловой системы: сектор не имеет имени файла. Как именно
BIOS взаимодействует с контроллером, этот вызов скрывает; при прямом
PIO CPU переносил бы слова сам, а при DMA контроллер передавал бы блок
после настройки. Ни DMA, ни собственный PIO-драйвер в действующей версии
журнала не реализованы.

\begin{center}\small
\begin{tabular}{|l|p{0.61\linewidth}|}\hline
Смещение & Формат одного сектора (512 байтов) \\ \hline
0--3 & ASCII \texttt{GR16} \\
4--7 & версия 1, 16 студентов, 4 работы, резерв 0 \\
8--71 & 64 оценки 0--10 или \texttt{FF} \\
72--73 & 16-битная сумма байтов 0--71, little-endian \\
74--511 & нули \\ \hline
\end{tabular}\end{center}
Новый формат \texttt{GR16} не смешиваем с историческим форматом
\texttt{GRD4} от x86-64-версии. Контрольная сумма обнаруживает некоторые
случайные повреждения, не защищает от преднамеренной подмены.
\texttt{save\_grades} зануляет 512-байтовый буфер, записывает поля по
фиксированным смещениям, считает сумму и вызывает \texttt{disk\_io}.
\texttt{load\_grades} проверяет диск, сигнатуру, версию, размеры,
сумму и все значения, только затем копирует таблицу в RAM. При ошибках
\texttt{EMPTY}, \texttt{NO\_DISK}, \texttt{BAD\_FORMAT},
\texttt{UNSUPPORTED\_VERSION}, \texttt{BAD\_CHECKSUM}, \texttt{BAD\_GRADE}
текущая таблица остаётся неизменной. Одна перезаписываемая копия сектора
не гарантирует атомарности при внезапной потере питания.

\paragraph{Машина в целом и переход к ОС.}
Плата и контроллеры соединяют CPU, RAM, накопитель, экран и клавиатуру.
В современном ПК для периферии используют USB, SATA и PCIe; NVMe ---
протокол хранения, а M.2 --- разъём/форм-фактор. Дискретный GPU
с собственной VRAM ориентирован на много подходящих параллельных операций;
наш текстовый VGA такой функциональности не реализует. HDD хранит данные
на вращающихся пластинах, SSD --- во Flash с контроллером. Их задержка
и пропускная способность различаются, хотя наш эмулируемый диск служит
лишь проверкой поведения, не скорости. BIOS смог помочь при запуске,
вводе и диске, но организации процессов, защиты памяти и файлов он
журналу не предоставил. Этим будет заниматься операционная система
в следующей лекции.

\paragraph{Самопроверка.}
Что произойдёт до \texttt{1000:0000}? Чем вызов \texttt{INT 16h}
отличается от IRQ1? Почему 0 --- оценка, а \texttt{FF} --- пропуск?
Почему \texttt{B800:0000} не равно адресу массива \texttt{grades}?
Что после перезапуска остаётся на дискете, а что на отдельном HDD?
Почему версия без гостевой ОС всё же может вызвать BIOS \texttt{INT 13h}?

Исходники действующего журнала: \path{demos/lecture04-journal/}.
Команда \texttt{make lecture4-journal-check} собирает шесть образов в Docker
и проверяет их запуск в QEMU; \texttt{make lecture4-journal-run}
показывает последнюю версию в браузере через noVNC по адресу
\url{http://localhost:6080/vnc.html}. Архивные 64-битные исходники
и прежняя презентация сохранены в \path{legacy/} для сравнения, но
не являются частью этого учебного маршрута.
\endgroup
